\chapter*{תקציר}
\addcontentsline{toc}{chapter}{תקציר}
\markboth{תקציר}{}

לפני מאתיים שנה, מידע נע במהירות של סוס. היום הוא נע במהירות האור, והוא כמעט אף פעם לא נח. הספר הזה שואל שאלה פשוטה לכאורה: מה קורה לבני אדם כשהם מחוברים תמיד?

כדי לענות עליה, הספר מתחיל רחוק מאוד מהמסך. הוא חוזר אל המברק של המאה ה-19, עובר דרך הרשת הצבאית שהפכה לאינטרנט, ומגיע אל המכשיר הקטן שבכיס, זה שבלע את המצלמה, המפה, הדואר והטלוויזיה. משם הוא עובר לנתונים: כ-5.87 מיליארד בני אדם השתמשו באינטרנט ב-2024 \parencite{itu2026data}, ומשתמש טיפוסי מבלה ברשת כ-6 שעות ו-40 דקות ביום \parencite{kemp2024}.

בהמשך, הספר בוחן את \textbf{כלכלת הקשב}, המנגנון שהופך כל גלילה לסחורה; את הקשר המורכב בין זמן מסך לרווחה נפשית, בעזרת מודל גולדילוקס פשוט, להמחשה בלבד, שמזמין לשאול לא רק "כמה מסך" אלא "כמה מכל דבר", ומזכיר שהאפקט שנמדד בפועל קטן; ואת הבית החכם, שבו הנוחות והפרטיות מושכות לכיוונים מנוגדים. הספר נחתם בארגז כלים מעשי לאיזון דיגיטלי ובמבט אל עידן הבינה המלאכותית.

אין כאן קריאה לזרוק את הטלפון לים. יש כאן הזמנה לעצור לרגע, להרים את הראש מהמסך, ולשאול מי בעצם מחזיק במי.

\keywords{זמן מסך, כלכלת הקשב, רווחה נפשית, השערת גולדילוקס, האינטרנט של הדברים, איזון דיגיטלי, \mbox{\en{Digital Well-being}}}
